\documentclass[11pt]{amsart}
\usepackage{amsmath,amssymb,amsthm}
\usepackage{graphicx}
\usepackage{booktabs}
\usepackage[margin=1in]{geometry}
\usepackage{hyperref}

\newtheorem{theorem}{Theorem}[section]
\newtheorem{proposition}[theorem]{Proposition}
\newtheorem{lemma}[theorem]{Lemma}
\theoremstyle{remark}
\newtheorem{remark}[theorem]{Remark}

\newcommand{\ve}{\varepsilon}
\newcommand{\R}{\mathbb{R}}

\title[Numerics for the high activation energy limit]{Computing the high activation energy limit\\ of an equidiffusional flame model}

\author{Juli\'an Fern\'andez Bonder}
\address{Departamento de Matem\'atica, FCEyN, Universidad de Buenos Aires, Argentina}

\date{DRAFT, \today}

\begin{abstract}
We study numerically the singular limit $\ve\to0$ of the system
$\Delta u-u_t=Yf_\ve(u)$, $\Delta Y-Y_t=Yf_\ve(u)$, $f_\ve(s)=\ve^{-2}f(s/\ve)$,
which models premixed flames with unit Lewis number in the limit of high
activation energy. When $(Y-u)/\ve\to w_0$, limit functions solve the free
boundary problem $u_t=\Delta u$ in $\{u>0\}$, $|\nabla u|=\sqrt{2M(w_0)}$ on
$\partial\{u>0\}$, so the limit depends on order-$\ve$ perturbations of the
data. We propose a finite volume method on a mesh that follows the flame
front (refining ahead of it and coarsening behind it), validate it against
exact travelling waves and the self-similar solution of Caffarelli and
V\'azquez, and observe first-order convergence in $\ve$ of the flame position
and speed in one space dimension and in radial symmetry ($N=2,3$). For
travelling waves we derive the first-order correction of the flame speed,
$c_\ve=c_0\big[1+\tfrac{\ve}{A}\big(1-\tfrac{3\sqrt3-2}{5}(1+w_0)\big)\big]+o(\ve)$
for $f(s)=(1-s)^+$, which changes sign at $w_0\approx0.564$, and confirm it
numerically.
\end{abstract}

\maketitle

% ---------------------------------------------------------------------------
\section{Introduction}

The system
\begin{equation}\label{eq:system}
\Delta u^\ve-u^\ve_t=Y^\ve f_\ve(u^\ve),\qquad
\Delta Y^\ve-Y^\ve_t=Y^\ve f_\ve(u^\ve),\qquad f_\ve(s)=\frac1{\ve^2}f\Big(\frac s\ve\Big),
\end{equation}
arises in the theory of equidiffusional premixed flames (Lewis number one)
\cite{BuckmasterLudford1982,ZeldovichFrankKamenetskii1938}: $Y^\ve\ge0$ is the
mass fraction of the reactant, $u^\ve=\lambda(T_f-T^\ve)$ a rescaled
temperature, and $\ve$ is essentially the inverse of the activation energy;
$f$ is Lipschitz, positive on $(-\infty,1)$ and zero on $[1,\infty)$. Since
$w^\ve=Y^\ve-u^\ve$ is caloric, \eqref{eq:system} reduces to
\begin{equation}\label{eq:Peps}\tag{$P_\ve$}
u_t=\Delta u-(u+w^\ve)f_\ve(u),\qquad w^\ve_t=\Delta w^\ve.
\end{equation}

In the case $u^\ve=Y^\ve$ the limit $\ve\to0$ was studied by Caffarelli and
V\'azquez \cite{CaffarelliVazquez1995}, by Caffarelli, Lederman and Wolanski
\cite{CaffarelliLedermanWolanski1997a,CaffarelliLedermanWolanski1997b} and, regarding uniqueness, by Lederman, V\'azquez and
Wolanski \cite{LedermanVazquezWolanski2001} and Petrosyan \cite{Petrosyan2002}. In
\cite{FernandezBonderWolanski2000} Wolanski and the author showed that if
$w^\ve/\ve\to w_0>-1$ then limits of $u^\ve$ solve
\begin{equation}\label{eq:P}\tag{$P$}
u_t=\Delta u\ \text{ in }\{u>0\},\qquad |\nabla u|=\sqrt{2M}\ \text{ on }\partial\{u>0\},\qquad
M=\int_{-w_0}^1(s+w_0)f(s)\,ds,
\end{equation}
in a pointwise and viscosity sense, and in \cite{FernandezBonderWolanski2004} that, under
geometric assumptions, the limit is unique and coincides with the classical
solution when one exists. Thus the high activation energy limit is
\emph{unstable}: perturbations of order $\ve$ of the initial data, which do
not change the limit of $u^\ve_0$, change the free boundary condition.

To the best of our knowledge [TO BE CONFIRMED WITH THE LITERATURE REVIEW] no
numerical study of the limit $\ve\to0$ in \eqref{eq:system} is available. The
difficulty is the reaction layer of width $O(\ve)$, which moves with the
flame and must be resolved: a uniform mesh needs $O(\ve^{-N})$ cells. In this
paper we
\begin{enumerate}
\item propose a finite volume method on a mesh that follows the flame,
with $O(\ve^{-1})$ cells in the layer and a geometric transition to a coarse
mesh elsewhere (Section~\ref{sec:method});
\item build exact reference solutions: travelling waves of \eqref{eq:Peps}
computed by a shooting method with a monotone energy, and the radial
self-similar solution of \eqref{eq:P}, which we write in terms of Kummer's
function (Section~\ref{sec:exact});
\item derive the first-order correction of the travelling wave speed
(Proposition~\ref{prop:speed});
\item observe first-order convergence in $\ve$ of front positions and
speeds, and display the instability of the limit with respect to $w_0$
(Section~\ref{sec:experiments}).
\end{enumerate}

Throughout, $f(s)=(1-s)^+$, for which
\begin{equation}\label{eq:M}
M(w_0)=\frac{(1+w_0)^3}{6}.
\end{equation}

% ---------------------------------------------------------------------------
\section{Exact solutions}\label{sec:exact}

\subsection{Travelling waves}
For \eqref{eq:P} in one dimension, $u=A\big(1-e^{-c_0(x-c_0t)}\big)^+$ is a
solution moving into the fresh region $\{u>0\}$ with speed
$c_0=\sqrt{2M}/A$, where $A=\lim_{x\to\infty}u$.

For \eqref{eq:Peps} with $w^\ve\equiv\ve w_0$ constant, travelling waves
$u=U(x-ct)$ satisfy $U''+cU'=(U+\ve w_0)f_\ve(U)$, $U(-\infty)=-\ve w_0$,
$U(+\infty)=A$. Ahead of the reaction zone ($U\ge\ve$) we have
$U=A-(A-\ve)e^{-c(\xi-\xi_1)}$. In the variables $U=\ve\varphi$,
$\xi-\xi_1=\ve z$, $\sigma=\ve c$,
\begin{equation}\label{eq:tw}
\varphi''+\sigma\varphi'=(\varphi+w_0)f(\varphi),\qquad
\varphi(0)=1,\quad \varphi'(0)=\sigma\Big(\frac A\ve-1\Big),\qquad \varphi(-\infty)=-w_0.
\end{equation}
Let $G(\varphi)=\int_{-w_0}^{\varphi}(s+w_0)f(s)\,ds$ (so $G(1)=M$) and
$E=\tfrac12\varphi'^2-G(\varphi)$. Along solutions,
$\frac{d}{dz}E=-\sigma\varphi'^2\le0$. Integrating \eqref{eq:tw} backwards from
$z=0$, if $\sigma$ is too large $E$ remains positive and $\varphi$ crosses
$-w_0$; if $\sigma$ is too small, $\varphi'$ vanishes before. This gives a
bisection for $\sigma$ (hence $c_\ve=\sigma/\ve$) to machine precision,
independent of the PDE solver.

\begin{proposition}\label{prop:speed}
Let $I=\int_{-w_0}^1\sqrt{2G(\varphi)}\,d\varphi$. Then
\[
c_\ve=c_0\Big[1+\frac\ve A\Big(1-\frac{I}{\sqrt{2M}}\Big)\Big]+o(\ve).
\]
For $f(s)=(1-s)^+$, $I/\sqrt{2M}=\kappa(1+w_0)$ with
$\kappa=\frac{3\sqrt3-2}{5}\approx0.6392$.
\end{proposition}

\begin{proof}[Sketch of proof]
Integrating $\frac{d}{dz}E=-\sigma\varphi'^2$ on $(-\infty,0)$ and using
$E(-\infty)=0$, $E(0)=\frac12\varphi'(0)^2-M$,
\[
\tfrac12\varphi'(0)^2=M-\sigma\int_{-\infty}^0\varphi'^2\,dz
=M-\sigma\int_{-w_0}^{1}\varphi'\,d\varphi .
\]
As $\sigma\to0$, $\varphi'\to\sqrt{2G(\varphi)}$ uniformly as a function of
$\varphi$, so the last integral is $I+o(1)$. Since
$\varphi'(0)=c(A-\ve)$ and $\sigma=\ve c$, we get
$c(A-\ve)=\sqrt{2M}-\ve cI/\sqrt{2M}+o(\ve)$, which is the claim. For
$f=(1-s)^+$ and $a=1+w_0$, $G=\frac a2t^2-\frac13t^3$ with $t=\varphi+w_0$, and
$I=\int_0^at\sqrt{a-\tfrac23t}\,dt=\tfrac94a^{5/2}\big[\tfrac4{15}-\tfrac23 3^{-3/2}+\tfrac25 3^{-5/2}\big]$;
with $\sqrt{2M}=a^{3/2}/\sqrt3$ this gives $I/\sqrt{2M}=\frac{3\sqrt3-2}5\,a$.
\end{proof}

In particular the first-order correction is positive for
$w_0<w_0^*=\frac5{3\sqrt3-2}-1\approx0.564$ and negative for $w_0>w_0^*$.
Table~\ref{tab:speedcoef} compares the prediction with the shooting method.

\begin{table}[h]
\centering
\begin{tabular}{rcccc}
\toprule
$w_0$ & Prop.~\ref{prop:speed} & $\ve=10^{-2}$ & $\ve=10^{-3}$ & $\ve=10^{-4}$\\
\midrule
$-0.5$ & $+0.138883$ & $+0.139876$ & $+0.138982$ & $+0.138893$\\
$0$    & $+0.208290$ & $+0.209515$ & $+0.208412$ & $+0.208303$\\
$1$    & $-0.454725$ & $-0.448230$ & $-0.454072$ & $-0.454660$\\
\bottomrule
\end{tabular}
\caption{$(c_\ve-c_0)/\ve$ from Proposition~\ref{prop:speed} and from the
shooting method ($A=1$).}\label{tab:speedcoef}
\end{table}

\subsection{The radial self-similar solution}
Caffarelli and V\'azquez \cite{CaffarelliVazquez1995} showed that \eqref{eq:P}
has solutions $V(x,t)=(T-t)^{1/2}h\big(|x|(T-t)^{-1/2}\big)$ with compact
support shrinking to a point at $t=T$; they are used in \cite[Prop.~4.3]{FernandezBonderWolanski2004}
to show that limit solutions vanish in finite time. Substituting,
$V$ is caloric in its positivity set if and only if
\begin{equation}\label{eq:h}
h''+\Big(\frac{N-1}r-\frac r2\Big)h'+\frac12h=0.
\end{equation}
(In \cite[(4.6)]{FernandezBonderWolanski2004} the sign of $r/2$ is misprinted.) The change of
variables $z=r^2/4$ turns \eqref{eq:h} into Kummer's equation
$zh_{zz}+(\frac N2-z)h_z+\frac12h=0$, so the solution regular at the origin is
\[
h(r)=\kappa_N\,{}_1F_1\big(-\tfrac12;\tfrac N2;\tfrac{r^2}4\big).
\]
Since $-1<-\frac12<0$, ${}_1F_1(-\frac12;\frac N2;\cdot)$ has exactly one
positive zero $z_N$; then $R=2\sqrt{z_N}$ and $\kappa_N$ is fixed by
$h'(R)=-\sqrt{2M}$. Numerically $R=1.8483$, $2.5139$, $3.0040$ for $N=1,2,3$.
The free boundary is the sphere of radius $R\sqrt{T-t}$.

% ---------------------------------------------------------------------------
\section{Numerical method}\label{sec:method}

We solve \eqref{eq:Peps} for $u$ and $w$ in planar ($N=1$) or radial
($N=2,3$) symmetry on $[0,L]$ with zero flux at both ends.

\emph{Space.} Vertex-centred finite volumes on a non-uniform mesh
$0=x_0<\dots<x_n=L$: with faces $x_{i\pm1/2}$ at the midpoints,
$|K_i|=(x_{i+1/2}^N-x_{i-1/2}^N)/N$ and
\[
(\Delta_hu)_i=\frac1{|K_i|}\Big(x_{i+1/2}^{N-1}\frac{u_{i+1}-u_i}{x_{i+1}-x_i}
-x_{i-1/2}^{N-1}\frac{u_i-u_{i-1}}{x_i-x_{i-1}}\Big).
\]
The scheme is conservative, exact on quadratics, and needs no special
treatment at $r=0$, where the weight $r^{N-1}$ vanishes.

\emph{Time.} The resulting stiff system is integrated with the variable
order, variable step BDF method of SciPy with the analytic sparse Jacobian
($\mathrm{rtol}=10^{-5}$, $\mathrm{atol}=10^{-4}\ve$).

\emph{Following the flame.} Every remeshing interval we
(i) locate the front $x_f$ as the point where $u=\ve$ (for $u\ge\ve$ there is
no reaction); (ii) define a band that covers the reaction layer ($5\ve$ on the
fresh side, $20\ve$ on the burnt side, where $u$ decays exponentially to
$-\ve w_0$) plus a margin of $15\ve$ on each side; (iii) build a new mesh
with spacing $h_f=\ve/8$ in the band and cells growing geometrically by a
factor $1.1$ away from it up to $h_c=0.01$; (iv) transfer $u,w$ to the new mesh
by monotone cubic (PCHIP) interpolation. The remeshing interval is chosen so
that the front, at its measured speed, moves at most half the margin.
Refinement ahead of the front and coarsening behind it happen in the same
step, and the number of nodes,
$O(L/h_c)+O(1/\ve\cdot\ve/h_f)+O(\log(h_c/h_f))$, does not depend on how far the
front travels.

\emph{Initial data.} Given the initial datum $u_0\ge0$ of \eqref{eq:P} we take
\[
u_0^\ve=u_0-\ve w_0\Big(1-\frac{u_0}\ve\Big)^+,\qquad Y_0^\ve=u_0^\ve+\ve w_0,
\]
so that $(Y_0^\ve-u_0^\ve)/\ve=w_0$ exactly, $Y_0^\ve=0$ in the burnt region
$\{u_0=0\}$, $Y_0^\ve\ge0$, and $u_0^\ve\to u_0$ uniformly with uniformly bounded
Lipschitz constants, as required in \cite{FernandezBonderWolanski2004}.

\begin{remark}
The cost is dominated by the number of time steps, which grows like
$\ve^{-1}$ to $\ve^{-3/2}$: on a mesh fixed between remeshings, the front
crosses $O(\ve^{-1})$ fine cells. A genuinely moving mesh (ALE or MMPDE
\cite{HuangRussell2011}), whose nodes travel with the front, would allow time
steps independent of $\ve$ in travelling-wave regimes.
\end{remark}

% ---------------------------------------------------------------------------
\section{Numerical experiments}\label{sec:experiments}

\subsection{Flame speed}
RESULTS_TW_TABLE

\subsection{Instability with respect to $w_0$}
RESULTS_TW_INSTABILITY

\subsection{Radial extinction}
RESULTS_RADIAL

% ---------------------------------------------------------------------------
\section{Conclusions and outlook}
OUTLOOK

\bibliographystyle{amsplain}
\bibliography{references}

\end{document}
